\chapter{Zadatak 2 --- što je točno zadano}
\label{ch:z2-zadatak}

\section{Izvorni tekst, doslovno}
\label{sec:z2-izvorni}

\begin{zadatak}{Zadatak 2 (Robustan QP)}
\skenzad[0.96]{z2-original}
\end{zadatak}

\begin{zadatak}{Zadatak 2 --- prijepis}
\textbf{Zadatak 2. (Robustan QP)} Za optimizacijski problem iz Zadatka 1
pretpostavimo da su koeficijenti $a_{ij}$, $i = 1,2$, $j = 1,2$ nesigurni u
sljedećem smislu. Jedino što znamo je da se njihove stvarne vrijednosti nalaze u
intervalu od $\pm 15\%$ od numeričkih vrijednosti zadanih u Zadatku 1: stvarna
vrijednost $a_{ij,\text{stvarno}}$ se nalazi u intervalu od $0.85 a_{ij}$ do
$1.15 a_{ij}$, gdje je $a_{ij}$ numerička vrijednost iz Zadatka 1. Riješite ovako
zadan robustan optimizacijski problem. Nacrtajte dozvoljen skup za niz slučajno
izabranih koeficijenata iz dozvoljenih intervala te na istoj slici označite gdje
se nalazi točka optimuma.
\end{zadatak}

\section{Što u zadatku \emph{ne} piše}
\label{sec:z2-nije-receno}

Zadatak 2 je najkraće formuliran od sva četiri, a traži najviše vlastitih odluka.

\begin{longtable}{@{}p{0.36\linewidth} p{0.60\linewidth}@{}}
\toprule
\textbf{Što nije rečeno} & \textbf{Kako je protumačeno i zašto} \\
\midrule
\endfirsthead
\toprule
\textbf{Što nije rečeno} & \textbf{Kako je protumačeno i zašto} \\
\midrule
\endhead
Što znači ,,riješiti robustan problem``. & Da se traži $x$ koji je dopustiv za
\emph{svaku} realizaciju koeficijenata iz intervala, i među takvima najbolji. To
je najgori slučaj (\emph{worst case}), a ne očekivana vrijednost ni neka
vjerojatnosna garancija --- zadatak ne spominje nikakvu razdiobu. \\[6pt]
Jesu li koeficijenti nezavisni. & Zadatak daje interval za svaki $a_{ij}$
zasebno, bez ikakve veze među njima. Uzima se dakle da svaki varira nezavisno,
pa je skup nesigurnosti \emph{pravokutnik} (kvadar), a ne npr.\ elipsoid ili
skup s ograničenim zbrojem odstupanja. \\[6pt]
Što se događa s $\pm15\%$ kad je nominalna vrijednost negativna. & Za
$a_{11} = -1$ interval ,,od $0.85a_{ij}$ do $1.15a_{ij}$`` doslovno znači od
$-0.85$ do $-1.15$, dakle u rastućem poretku $[-1.15,\, -0.85]$. Zadatak
navodi granice u tom redoslijedu, pa ih treba preurediti; vidi
zamku u \S\ref{sec:z2-intervali}. \\[6pt]
Je li i funkcija cilja nesigurna. & Nije --- nesigurnost je izričito vezana samo
uz $a_{ij}$. Matrice $Q$ i $c$ ostaju iste kao u Zadatku 1, kao i desne strane
$-10$ i $3$. \\[6pt]
Koliko realizacija nacrtati. & Kaže se samo ,,niz slučajno izabranih``. U
seminaru je uzeto $100$, uz \texttt{rng(42)} radi ponovljivosti. \\[6pt]
Treba li usporediti s nominalnim rješenjem. & Ne traži se izrijekom, ali je to
jedini način da se pokaže \emph{čemu} robusnost služi i koliko košta. U seminaru
je dodana usporedba i brojanje kršenja po realizacijama. \\
\bottomrule
\end{longtable}

\begin{zamka}[Zamka: ,,dozvoljen skup za niz slučajnih koeficijenata``]
Zadatak traži da se nacrta dopustivi skup \emph{za realizacije}, ne skup
nesigurnosti. To su dvije različite stvari u dva različita prostora:
\begin{itemize}[nosep,leftmargin=*]
  \item \textbf{skup nesigurnosti} $\mathcal{P}_i$ živi u prostoru
        \emph{koeficijenata} $(a_{i1}, a_{i2})$ --- to je pravokutnik;
  \item \textbf{dopustivi skup} živi u prostoru \emph{projektnih varijabli}
        $(x_1, x_2)$ --- i za svaku realizaciju koeficijenata izgleda malo
        drukčije.
\end{itemize}
Slika koju zadatak traži je u prostoru $(x_1, x_2)$: mnogo blago različitih
rubnih pravaca, po jedan par za svaku realizaciju.
\end{zamka}

\section{Karta poglavlja}
\label{sec:z2-karta}

\begin{longtable}{@{}p{0.48\linewidth} p{0.48\linewidth}@{}}
\toprule
\textbf{Korak} & \textbf{Gdje je obrađen} \\
\midrule
\endhead
Skup nesigurnosti i robusna formulacija & Poglavlje~\ref{ch:z2-nesigurnost} \\
Svođenje na konačno mnogo nejednakosti & Poglavlje~\ref{ch:z2-vrhovi} \\
Analitičko rješenje, KKT provjera, MATLAB & Poglavlje~\ref{ch:z2-rjesenje} \\
Slučajne realizacije, slika, usporedba & Poglavlje~\ref{ch:z2-realizacije} \\
Pitanja, zamke i modelski odgovori & Poglavlje~\ref{ch:z2-obrana} \\
\bottomrule
\end{longtable}
